% =========================================================
% CUSTOM COMMANDS AND SHORTCUTS
% =========================================================
% This file defines reusable LaTeX commands that simplify
% common tasks throughout the document.
%
% HOW TO USE: Call these commands in your chapter files
% to maintain consistency and save typing effort.
% =========================================================

% ---------------------------------------------------------
% FONT SIZE SHORTCUTS
% ---------------------------------------------------------
% These commands provide named font sizes matching the
% Anna University guidelines. Use them on the title page,
% certificates, and headings.
%
% USAGE: {\TitleSize This text will be 18pt}
% ---------------------------------------------------------
\newcommand{\TitleSize}{\fontsize{18}{27}\selectfont}
\newcommand{\SubTitleSize}{\fontsize{16}{24}\selectfont}
\newcommand{\NormalSize}{\fontsize{14}{21}\selectfont}
\newcommand{\SmallSize}{\fontsize{12}{18}\selectfont}

% ---------------------------------------------------------
% FIGURE INSERTION COMMAND
% ---------------------------------------------------------
% Inserts a figure with proper formatting in one command.
%
% USAGE:
%   \insertfigure{width=0.8\textwidth}{filename}{Caption text}{fig:label}
%
% PARAMETERS:
%   #1 = width specification (e.g., width=0.8\textwidth)
%   #2 = image filename (relative to figures/ folder)
%   #3 = caption text
%   #4 = label for cross-referencing (use \ref{fig:label})
%
% EXAMPLE:
%   \insertfigure{width=0.8\textwidth}{use_case_diagram}{Use Case Diagram of the System}{fig:usecase}
%   As shown in Figure~\ref{fig:usecase}, the system...
% ---------------------------------------------------------
\newcommand{\insertfigure}[4]{%
  \begin{figure}[htbp]
    \centering
    \includegraphics[#1]{#2}
    \caption{#3}
    \label{#4}
  \end{figure}
}

% ---------------------------------------------------------
% TABLE INSERTION HELPER
% ---------------------------------------------------------
% For creating a simple table with a caption above it.
% For complex tables, write the tabular environment directly.
%
% USAGE:
%   \begin{table}[htbp]
%     \centering
%     \caption{Your table caption}
%     \label{tab:yourlabel}
%     \begin{tabular}{|c|c|c|}
%       \hline
%       Header 1 & Header 2 & Header 3 \\
%       \hline
%       Data 1   & Data 2   & Data 3   \\
%       \hline
%     \end{tabular}
%   \end{table}
%
% Remember: Table captions go ABOVE the table (per AU guidelines).
% Figure captions go BELOW the figure.
% ---------------------------------------------------------

% ---------------------------------------------------------
% CROSS-REFERENCE SHORTCUTS
% ---------------------------------------------------------
% These commands ensure consistent formatting when referring
% to figures, tables, chapters, and equations in the text.
%
% USAGE:
%   As shown in \figref{fig:usecase}, the system...
%   The values are listed in \tabref{tab:results}.
%   This is discussed in \chapref{chap:intro}.
%   From \eqref{eq:formula}, we can derive...
% ---------------------------------------------------------
\newcommand{\figref}[1]{Figure~\ref{#1}}
\newcommand{\tabref}[1]{Table~\ref{#1}}
\newcommand{\chapref}[1]{Chapter~\ref{#1}}

% ---------------------------------------------------------
% SIGNATURE BLOCK COMMAND
% ---------------------------------------------------------
% Creates a signature block used on certificate pages.
%
% USAGE:
%   \signatureblock{Dr. Name}{DESIGNATION}{Department}{Address}
%
% PARAMETERS:
%   #1 = Name
%   #2 = Designation/Role (e.g., SUPERVISOR, HEAD OF THE DEPARTMENT)
%   #3 = Department
%   #4 = Address
% ---------------------------------------------------------
\newcommand{\signatureblock}[4]{%
  \begin{minipage}[t]{0.45\textwidth}
    \vspace{1.5cm}
    \textbf{SIGNATURE}\\[0.5cm]
    \textbf{#1}\\
    \textbf{#2}\\[0.3cm]
    \textbf{#3}\\
    #4
  \end{minipage}%
}

% ---------------------------------------------------------
% BLANK PAGE COMMAND
% ---------------------------------------------------------
% Inserts a completely blank page (no headers, footers, or
% page numbers). Useful for the back of the title page in
% double-sided printing.
%
% USAGE: \blankpage
% ---------------------------------------------------------
\newcommand{\blankpage}{%
  \clearpage
  \thispagestyle{empty}
  \mbox{}
  \clearpage
}
